\thispagestyle{plain}
\ctsection{Præfatio}

\begin{multicols}{2}
Cum eorum, qui humaniores litteras repetebant, domesticæ scholæ præfectus fuissem, id propositum habebam maxime, et eo industriam omnem convertendam existimabam, ut P. Juvencii, harum litterarum hodie inter nos facile Principis, sapienti monito obtemperarem. Primus, ait ille, Magistri labor erit, Latinum stilum illis optimum demonstrare, et ut assequantur, adjuvare. Ergo putare mecum cœpi, quæ adjumenta offerrem præsentissima, queis viis Adolescentes non minus sciendi aviditate, quam ingenii præstantia, ad omnem hujus litteraturæ laudem paratissimos, ad Latinæ Linguæ Proprietatem, Elegantiam, Copiam quam celerrime commodissimeque perducerem.

\textit{Proprietatem} ut condiscerent, hoc est, e barbarie vulgaris sermonis sordibusque emergerent, præter expositos barbararum vocum Indices multiplices, si quam vocem suspectam haberent, \textit{Thesaurum Latinitatis}, seu \textit{Forum Romanum} consulerent, jussi. Ex eodem præcipuorum verborum syntaxin usumque Latinitatis proprium ipse exscripsi: quo labore perfunctus admoneor a Collega, dudum id a P. Pomæio in suis Flosculis præstitum. Etsi dolerem actum me egisse, deque libello tam pervulgato ignorasse puderet, in eamdem tamen cum viro egregio cogitationem, methodum ac prope verba incidisse mihi gratulabar. Hos ergo flores transcribendos dedi meis; seposui quot diebus quartam horulæ partem, qua suscitandos pertentarem singulos, constitutumque ejus diei pensum de memoria repeterem; si qua adhærerent, explicabam; difficiliora paradigmatis illustrabam; et fiebat, ut quotidianæ scriptionis argumento fere aliquid inserere possent. Atque hisce scriptionis, auditionis, exercitationis industriis jam intra bimestre Latinum quid velut in herba efflorescere, non sine jucundo animi sensu animadverti.

Ut vero ad \textit{Elegantiam} quoque perpolirentur, Tursellini particulas singulis, Schorum illis aut Doletum, Hadrianum aliis aut Goclenium, etc., ut copia erat, divisi, monuique, ne legerent duntaxat, sed lecta in chartam compendio referrent, quo memoriæ commendarent altius, quæ ab lectione sæpe, ut fit, minus intenta facile devolatura erant.

Ut \textit{Copiam} cum elegantia combiberent, Manutium a Buchlero auctum, et Schönslederum præbui, quos altera quavis periodo inspicerent, quorum opibus ad ornatum varietatemque nec timide et liberaliter tantisper uterentur, dum usu non admodum diutino in jus suum assererent quæ aliena fuerant, abjectisque centonibus, de suo liberum quid et nativum contexerent; cogitarent autem præterea: hæc incipientium subsidia ad proritandum duntaxat, Latinoque aliquo sapore imbuendum Tironum palatum excogitata; cæterum germanæ Latinitatis spem omnem in assidua Ciceronis, bonorumque ex aurea ætate Auctorum lectione collocandam; hos tanta cura pervolvendos velut si nihil peractum paratumque accepissent a majoribus, sed unis ipsis eruenda forent, quæ ex his venis tantis conatibus egessere, quidquid hactenus fuit Phraseologorum.

Hæc agenti cum illud incommodum videbatur, quod tot libellorum, queis Discipulorum pluteos onerarem, turba mihi undeunde corradenda esset, tum quod iisdem, dum hinc Adverbia, Epitheta inde, etc., fastidiose commendicant, plurimum utilissimi temporis volvendo revolvendoque inaniter elaberetur. Hoc tædio ut levaremur utrinque, optare cœpi, esset aliquis, qui sparsas has operas, eademque apud diversos repetita in unum quoddam corpus conflandi consilium caperet. Tot jam exstare Criticorum, Grammaticorum, Phraseologorum Indices, Enchiridia, Gazophylacia, Thesauros, Elegantias observavi, ut felicem se reputare posset is, quem tantam, tam paratam materiam in unum volumen certo ordine ac methodo ædificandi animus incesseret; feliciores vero fore Discipulos, qui in magno viæ compendio progressiones tamen celeriores facerent, nec jam per rivulos concursare deberent, quos in unum fontem Vir adeo beneficus corrivasset.

Externis vero nostris scholasticis quam ejusmodi opera non opportuna modo, sed necessaria foret omnino! Ciceronem, Ciceronem ex alto suggestus Juvenibus inclamamus magistri, in hunc velut regulam omnisque Latinitatis exemplum respicere, hunc legere, imitari et scribere jubemus, sed iis juvenibus, quorum parti maximæ, præter fragmenta quædam suæ classis scholastico libro (quamquam et hunc quam pauci emant) inserta nullus domi Cicero ; plurimis Terentii, Livii, Sallustii, etc. ipsa nomina peregrina, plerique, quid Nizzolius aut Forum Romanum hominis vel rei sint, perinde ut maragnonius quispiam ignorent. Et hos stimulis fodiemus? hos Latinos Ciceronianos effici volumus? qua autem ope, ad quem lydium lapidem suspectæ fidei voces explorent? barbaras secernant, melioribus commutent? unde colores, verasque pro adulterinis pigmentis elegantias petant? unde ad copiolam aliquam parandam vel Synonymiam exsugant? Huc certe flagitii vidi devenisse nonnullos, ut Epistolam, chriam, orationem scripturi \textit{in Synonymis}, \textit{Gradibus ad Parnassum}, etc. habitarent, qua re quid horribilius et ad corrumpendum stilum aptius excogitari possit? quod si opulentiores quidam filii familias circumfluant etiam natentque Ciceronibus, Terentiis, etc. quam pauci iis urantur igniculis, ut formulas inde, phrases, etc. in adversaria excerperent, justam animi contentionem ad exhaurienda principiorum ardua afferrent?

Eædem atque his majores difficultates, cum Poeticæ fores obsiderent, fuere qui \textit{Flavissas}, \textit{Synonyma}, \textit{Gradus ad Parnassum}, etc., præclara adeo adjumenta elaborarent, ut brevi tempore Poeticum aliquem colorem inde caloremque trahere multos experiamur. Quid est igitur, quod nemo dum ad Oratoriam gradientibus de tali adminiculo de gradibus ad rostra gratificatus est? quid? quod optarunt hoc multi, nemo optantibus morem gessit? quid autem (aiebam), quin tute hanc spartam occupas, et lexicon aliquod Phraseologicum, aut quodvis demum nomen habeat, vel ipse conficis, vel si per temporis angustias excludaris, typum denique et summas lineas designas, quas tali provocatione commoti alii depingant denique et absolvant.

Consilii dubium duæ res magnopere retrahebant, magnitudo laboris et ignobilitas. Infida utebar parumque firma valetudine. Operosæ scholæ quotidianæ vacui utilisque temporis relinquebant minimum; et informaveram nonnulla, quorum Auctor minimum dici possem, cum e præsenti opere præter diligentis scribæ nomen nihil ad me rediturum cernerem. Titubantem Collega meus, externorum Rhetorum Magister, Vir et morum et Latinitatis purissimæ candore longe suavissimus, P. Gabriel Wimmert, iteratis identidem hortatibus in proposito retinuit, tantoque promptiorem habuit, quod ejusmodi, ut ego sum, hominum genus, ab omni fumo alienissimum esse, idque pro honestissimo ducere debet, si, magno quidem suo impendio, quæstu minimo, plurimorum commodis servire possit.

Re deliberata Phraseologos, quos nancisci poteram, omnes tres in classes partior; horum alii \textit{Proprietatem} docerent, ut Niessii, Cellarii, Vossii, Vorstii, Borichii Commentarioli; alii \textit{Elegantiam} præstarent, ut Hadrianus, Valla, Doletus, Tursellinus, Schorus, Fabri, etc., \textit{Copiam} denique nonnulli suppeditarent, quorum chorum Manutius duceret et Schönslederus. Horum, inquam, operas omnes in unum Corpus cogendas, eximenda, quæ apud singulos præcipua essent, atque sub unum conspectum objicienda constitui, ea tamen adjecta conditione, ut ne voluminis moles ita excresceret, quin ab juvenibus in scholam ad statas scriptiones comportari posset. Hanc vero, quam vides, methodum mihi descripsi:

I. Nizzolium equidem pro solo operis stravi, ejusdemque ordinem alphabeticum sum secutus, quin tamen iisdem me carceribus includerem, iisdem religionibus obligarem; neque enim in ea sum hæresi, ut vel Ciceronem de omnibus argumentis disputasse, aut non iisdem cum Vitruvio, Varrone, Livio, Cæsare, etc., si de iisdem rebus loqui voluisset, vocabulis usurum fuisse, existimem. Quare aliorum quoque ex aurea ætate Auctorum voces adjeci, asterisco notatas, ut ab Ciceronianis secernantur, non quod fugiendæ essent. Totus enim in ea sum sententia, juvenes intra hujus ætatis limites constringendos, et a Plinianis, Senecianis, etc. omnino abstinendos. Intelligent autem, quam hic vocem non repererint, nullum habere idoneum Auctorem: atque adeo meliore commutandam.

II. Vocabula nihilominus \textit{sacra et religiosa} suis locis characterum forma diversa intexui, eademque e Pontano, quo nemo melius christiane simul et Latine loqui docet, Tympii opera collectis observationibus illustravi.

III. Per illud Cf. \textit{confer}, passim, qua eundum esset, præsignavi.

IV. Synonyma, ut in Nizzolio sunt, simpliciter adscripsi.

V. Contraria signo )( notata, ut quæ se ipsis offerrent, non erat, ut multa vestigatione conquirerem.

VI. Vocum, quas fere pro idem significantibus usurpat vulgus, elegantiores vero haud paulum notionibus inter se differre observant, \textit{Differentias}, e Laurentio Valla dedi aliquas, ac revera cœpi potius, et quid voluissem ostendi.

VII. Suam vocibus singulis syntaxim seu \textit{Usum} peculiarem, sententiis formulisque lectissimis e Nizzolio, Pomæio, aliisque multo discrimine decerptis attexui, quarum iterata lectio cum alias utilitates afferret, tum si quampiam atque ab his alienam construendi formulam adhibuerint Discipuli, fere in Idiotismum se impactos admoneantur.

VIII. Vulgatiores barbarismos seu Idiotismos Syntaxi Latinæ locis quibusdam e Schoro, etc. subjeci; ut venia sit ingenue confitenti, ubi minimum oportebat, præparcus fui; nec enim solum demonstranda erat via, qua incederent; sed infames quotidianis erroribus scopuli detegendi, ut suæ barbariei identidem admonerentur juvenes.

IX. Ad Copiam seu Phraseologiam quod attinet, nempe operis partem principem, ac discentibus desideratissimam, primæ Editiones nonnisi collectas Roterodami, Manutii, quique omnibus præstat, Schönslederi Phrases, paulum ad Juvenum captum deflexas habuere. In præsens opera cujusdam Collegæ mei, et tædii vere decumani labore, accessere Liviana, Cæsareana, Corneliana, Sallustiana, cum Indice diffusiore vocum barbararum, alioque Tironibus admodum desiderato Gallico-Latino; Opusque evasit ejusmodi, ut detrahi nihil, addi perparum possit. Ab viri modestia obtinere non potui, ut suum de studiosa Juventute tam præclare meriti nomen adscribi pateretur.

Habetis hic humaniorum litterarum Professores, quo consilio primum cœptum sit Opus, et quantum ad ejus absolutionem recens accesserit. Si præterquam quid præstiterimus, quid insuper voluerimus, apud vos existimabitis, laboribus nostris, vestro vestrorumque commodo benevoli utemini, fruemini.

\end{multicols}